\chapter{Methodology}
\label{ch:methodology}

\section{System overview}
\label{sec:overview}

The system is built as a five-layer stack (Figure~\ref{fig:architecture}) around
one governing idea: every claim an agent makes has to be traceable back to
something it was actually given. A document goes in, an extraction agent turns it
into a validated scenario object, four agents interact under an explicit state
machine, and every step is written to a JSON log that later analysis can read
without ever re-invoking the model. The Python backend runs to 33 modules and
roughly 4,900 lines, backed by a 92-case test suite; the API and frontend add a
further 374 and 3,800 lines on top of that.

\begin{figure}[htbp]
\centering
\resizebox{\textwidth}{!}{%
\begin{tikzpicture}[font=\footnotesize,
    badge/.style={circle,draw=f21ink,fill=white,text=f21ink,
                  font=\bfseries\footnotesize,minimum size=9mm,inner sep=0pt},
    layer/.style={draw=f21ink,line width=0.7pt,rounded corners=3pt,
                  text width=10.8cm,minimum height=1.15cm,align=left,
                  fill=f21mist,inner sep=7pt},
    gaplayer/.style={draw=f21grey,dashed,line width=0.7pt,rounded corners=3pt,
                text width=10.8cm,minimum height=0.95cm,align=left,
                fill=white,inner sep=7pt},
    tag/.style={f21grey,font=\scriptsize\itshape}]

  \node[badge] (b1) at (0,0) {L1};
  \node[layer,right=6mm of b1,anchor=west] (l1)
    {\textbf{Presentation} --- React/TypeScript client: upload, live streamed\\
     negotiation, case browser, analytics, risk screen \hfill\tag{browser}};

  \node[badge,below=8mm of b1] (b2) {L2.5};
  \node[layer,right=6mm of b2,anchor=west] (l2)
    {\textbf{API} --- FastAPI: extraction endpoint, server-sent-event\\
     negotiation stream, risk assessment, GPU-tunnel status and reconnect};

  \node[badge,below=8mm of b2] (b3) {L2};
  \node[layer,right=6mm of b3,anchor=west] (l3)
    {\textbf{Orchestration} --- LangGraph \texttt{StateGraph}: deterministic\\
     turn order, conditional routing, per-super-step streaming \hfill\tag{process}};

  \node[badge,below=8mm of b3] (b4) {L3};
  \node[layer,right=6mm of b4,anchor=west] (l4)
    {\textbf{Agents} --- Contracting Authority, Aggrieved Bidder, Court,\\
     Summary, Extraction (Llama 3.1 8B via Ollama; local CPU or\\
     tunnelled A10G GPU) \hfill\tag{model}};

  \node[badge,below=8mm of b4] (b5) {L4};
  \node[gaplayer,right=6mm of b5,anchor=west] (l5)
    {\textbf{Knowledge / RAG} --- vector retrieval over legislation and\\
     case law \emph{--- deliberately not built; see \S\ref{sec:alternatives}}};

  \node[badge,below=8mm of b5] (b6) {L5};
  \node[layer,right=6mm of b6,anchor=west] (l6)
    {\textbf{Persistence} --- structured JSON logs per run; batch summaries\\
     with provenance and compliance counters};

  \draw[->,f21ink] (b1) -- (b2);
  \draw[->,f21ink] (b2) -- (b3);
  \draw[->,f21ink] (b3) -- (b4);
  \draw[->,f21grey,dashed] (b4) -- (b5);
  \draw[->,f21ink] (b5) -- (b6);
\end{tikzpicture}}
\caption[System architecture]{The five-layer architecture, read top to bottom as
the path a request takes. The dashed layer is reported because it was scoped
out, not because it exists: the extraction agent is a single-pass synthesis
step and must not be mistaken for retrieval.}
\label{fig:architecture}
\end{figure}

\section{Grounding the agents in the client's own material}
\label{sec:grounding}

An early version of this system used interest lists I had written myself, which
made the agents believable but gave nobody a way to challenge them. Those were
replaced with the taxonomies already sitting in Fusion21's Doc 1, so that every
agent's stated interests trace back to a named category in a document the client
wrote, rather than to my own guess at what a procurement officer cares about. Doc
1 sets out seven categories for the contracting authority, six for the supplier,
and six for the judiciary; Table~\ref{tab:agent-config} summarises the resulting
configuration.

Doc 1 also supplies per-case behavioural modifiers its ``practical
considerations'' tables which vary by dispute rather than sitting fixed to a
role. These are modelled as optional \texttt{CAProfile} and \texttt{BidderProfile}
objects attached to a scenario and rendered into the system prompt at
construction time. Both default to unset and render no text at all, so that every
result obtained before profiles existed stays directly comparable to later ones.
That decision carries a structural consequence: agents can no longer be
module-level singletons, since their system prompts now depend on the scenario in
front of them, so they get built per dispute and bound into the graph nodes at
compile time instead.

\begin{table}[htbp]
\centering
\footnotesize
\caption[Agent configuration]{Agent configuration. Temperatures were set by role
function rather than tuned: the adjudicator and summariser are held low because
consistency matters more than variety, the challenger is highest because a
negotiation in which nobody proposes anything new stalls.}
\label{tab:agent-config}
\begin{tabular}{@{}llcp{5.4cm}@{}}
\toprule
\textbf{Agent} & \textbf{Interest source} & \textbf{Temp.} & \textbf{Structured output} \\
\midrule
Contracting Authority & Doc 1, 7 CA categories & 0.3 &
  \texttt{PreNegotiationStatement}, \texttt{RoundResponse}, \texttt{WinStatement} \\
Aggrieved Bidder & Doc 1, 6 supplier categories & 0.4 &
  \texttt{PreNegotiationStatement}, \texttt{RoundResponse}, \texttt{WinStatement} \\
Court / Judge & Doc 1, 6 judiciary categories & 0.2 &
  \texttt{ComplianceAssessment} \\
Summary (non-negotiating) & --- & 0.2 &
  \texttt{NegotiationSummary} \\
Extraction & --- & 0.1 &
  \texttt{DisputeScenario} \\
\bottomrule
\end{tabular}
\end{table}

The fourth agent was added on my supervisor's suggestion and takes no part in the
negotiation itself. It reads the finished transcript and writes a plain-English
account of what happened, on the reasoning that a system whose output only a
procurement lawyer can parse has not really lowered anybody's costs.
Section~\ref{sec:readability} tests that claim against a standard readability
measure rather than leaving it as an assertion, which turns out to matter: the
summaries read fluently enough that the shortfall is invisible without measuring
it.

\section{The Court agent: conditional verification}
\label{sec:court-design}

The Court agent assesses process compliance only, which mirrors how UK courts
actually review procurement decisions: the live question is whether the authority
followed a lawful, rational and procedurally correct process, not whether the
better bidder happened to win. Getting the agent to behave took four revisions,
each one driven by a measured failure rather than a hunch
(Table~\ref{tab:court-versions}).

\begin{table}[htbp]
\centering
\footnotesize
\setlength{\tabcolsep}{4pt}
\renewcommand{\arraystretch}{1.25}
\caption[Court prompt revisions]{Four revisions of the Court agent's prompt. Each
row's failure mode is what motivated the next row.}
\label{tab:court-versions}
\begin{tabular}{@{}p{1.4cm}p{3.6cm}p{5.4cm}@{}}
\toprule
\textbf{Version} & \textbf{Instruction} & \textbf{Observed behaviour} \\
\midrule
V1 & Judgement-based, no explicit check &
  Manifest-error findings tracked whether the CA conceded in dialogue, not the
  facts \\
V2 & Always verify any calculation &
  Fixed the numeric cases, but invented point systems on qualitative disputes
  that had none \\
V3 & Conditional: compute if a real formula exists, otherwise qualitative
  review &
  Restored non-fabricating judgement while keeping numeric accuracy \\
V4 & V3 with the judiciary interest block replaced by Doc 1's six categories &
  Never deadlocks on any case; see \S\ref{sec:ablation} for whether that is
  good news \\
\bottomrule
\end{tabular}
\end{table}

V4 is derived from V3 through a single verified string substitution of one named
block, and the module raises an error at import if that anchor ever stops
matching. The two prompts therefore differ in exactly one variable, which is what
turns the ablation in Section~\ref{sec:ablation} into a controlled comparison
rather than a before-and-after anecdote.

The decision procedure itself is shown in Figure~\ref{fig:court-flow}. Its
least obvious part is the middle branch. A scenario that states a formula and a
correction, but not the original inputs, is not actually computable, and an
earlier version of the prompt let the model paper over that gap by assuming the
missing figures. What that produced was arithmetic that looked internally valid
and was evidentially worthless. Step 1's gate was tightened accordingly, to
require the formula \emph{and} every input value it needs, with an explicit note
that a stated correction describes an outcome rather than an input, backed by an
outright ban on writing down an unstated number even when it is hedged as
illustrative or assumed. The hedge is not a safeguard here; it is the exact shape
the failure took.

\begin{figure}[htbp]
\centering
\resizebox{\textwidth}{!}{%
\begin{tikzpicture}[
    font=\footnotesize,
    >=Latex,
    node distance=8mm and 6mm,
    start/.style={draw=f21ink,rounded corners=3pt,align=center,inner sep=6pt,
                  fill=f21mist,line width=0.8pt,text width=5.0cm},
    decision/.style={draw=f21ink,rounded corners=3pt,align=center,inner sep=7pt,
                     fill=white,line width=1.3pt,text width=5.6cm},
    outcome/.style={draw=f21green,line width=1pt,rounded corners=3pt,align=center,
                    inner sep=6pt,fill=f21green!8,text width=3.9cm},
    warn/.style={draw=f21amber,line width=1pt,rounded corners=3pt,align=center,
                 inner sep=6pt,fill=f21amber!8,text width=4.6cm},
    lbl/.style={font=\scriptsize\itshape,fill=white,inner sep=1.5pt}
  ]

  \node[start] (in) {Round transcript + scenario};
  \node[decision, below=of in] (q1)
      {\textbf{Step 1}\\[2pt] Does the scenario state a formula \emph{and}
       every input value that formula needs?};

  \node[outcome, below left=14mm and 2mm of q1] (a)
      {\textbf{Step 2A --- Verify}\\ Recompute exactly, show working.\\
       Mismatch $\Rightarrow$ manifest error};
  \node[outcome, below right=14mm and 2mm of q1] (b)
      {\textbf{Step 2B --- Review}\\ Qualitative rational-basis review.\\
       No invented numbers};
  \node[warn, below=24mm of q1] (c)
      {\textbf{Formula stated, inputs missing}\\
       Do not assume, estimate, or ``illustrate'' the gap};

  \draw[->] (in) -- (q1);
  \draw[->] (q1.west) -- ++(-4mm,0) |- node[lbl,pos=0.7] {both present} (a.north);
  \draw[->] (q1.east) -- ++(4mm,0) |- node[lbl,pos=0.7] {neither} (b.north);
  \draw[->] (q1.south) -- node[lbl,right] {formula only, inputs missing} (c.north);
  \draw[->,f21amber,line width=1.1pt] (c.south) -- ++(0,-4mm) -| node[lbl,below,pos=0.75]
      {falls back to 2B} (b.south);
\end{tikzpicture}}
\caption[Court agent decision procedure]{The Court agent's gating procedure. The
amber path is the one that had to be added after the Woods v Milton Keynes case
exposed it: a formula plus a stated correction is not a computable scenario, and
treating it as one produced fabricated inputs.}
\label{fig:court-flow}
\end{figure}

The same discipline applies to legal citation. The agent may only cite a
provision if it appears in the scenario, in a party's statement that round, or in
the grounding material handed to it in its own system prompt. A correct general
point with no citation attached is explicitly preferred over a specific-sounding
invented one, on the view that a fabricated regulation number belongs to exactly
the same category of error as a fabricated sub-score. Section~\ref{sec:citation}
measures whether that instruction actually holds up.

\section{Orchestration}
\label{sec:orchestration}

Turn order runs through a LangGraph \texttt{StateGraph} (Figure~\ref{fig:graph})
rather than a conversational free-for-all. Both parties set out their interests,
goals and BATNA before any negotiating starts; then the authority and the bidder
alternate, with the Court checking compliance after every round. One conditional
edge decides whether the loop continues.

\begin{figure}[htbp]
\centering
\resizebox{\textwidth}{!}{%
\begin{tikzpicture}[font=\footnotesize,>=Latex,
  n/.style={draw=f21ink,rounded corners=2pt,minimum height=8mm,inner xsep=6pt,
            fill=white,line width=0.7pt},
  c/.style={draw=f21green,rounded corners=2pt,minimum height=8mm,inner xsep=6pt,
            fill=f21green!8,line width=0.9pt}]
  \node[n] (pre) at (0,0) {pre\_negotiation};
  \node[n,right=9mm of pre] (ca) {ca\_round};
  \node[n,right=9mm of ca] (bid) {bidder\_round};
  \node[c,right=9mm of bid] (crt) {court\_check};
  \node[n,right=11mm of crt] (win) {win\_statements};
  \node[n,below=9mm of win] (sum) {summary};
  \node[n,right=9mm of sum] (end) {\textsc{end}};

  \draw[->] (pre) -- (ca);
  \draw[->] (ca) -- (bid);
  \draw[->] (bid) -- (crt);
  \draw[->] (crt) -- node[above,font=\scriptsize,align=center]
      {resolved\\ or round $\geq$ max} (win);
  \draw[->] (win) -- (sum);
  \draw[->] (sum) -- (end);
  \draw[->,f21green] (crt.south) -- ++(0,-0.9) -| node[below,pos=0.25,font=\scriptsize]
      {\textcolor{f21green}{continue}} (ca.south);
\end{tikzpicture}}
\caption[Negotiation state graph]{The negotiation graph. Everything hinges on
\texttt{court\_check}: it is the only node that can set \texttt{resolved}, and
therefore the only escape from the loop other than exhausting the round limit.
Section~\ref{sec:baselines} removes it and measures the consequence.}
\label{fig:graph}
\end{figure}

Resolution is decided by the Court's \texttt{recommended\_action}, drawn from a
closed vocabulary: continue negotiation, re-evaluation, no remedy with the
decision standing, or damages. Anything other than ``continue'' resolves the
dispute; exhausting the round limit without resolution is recorded as deadlock and
escalation to formal proceedings. Constraining this vocabulary is a legal
requirement rather than an implementation convenience (Doc 3 is explicit that
award-stage disputes cannot be settled by splitting the difference), and the
constraint is enforced rather than merely documented: the downstream
recommendation module raises an error if a batch contains any outcome outside the
set, which was checked by deliberately injecting the string ``settle for 50\% of
contract value'' during testing.

The orchestrator exposes both a synchronous \texttt{run()} used by batch scripts
and a \texttt{stream()} generator used by the API. One detail is worth recording
here because it caused a real bug in practice: the streamed per-node updates are
partial, and \texttt{bidder\_round} does not carry \texttt{round\_number} in its
update the way \texttt{ca\_round} does so any consumer has to read that field from
merged state rather than from the raw update itself.

\section{From documents to scenarios}
\label{sec:extraction}

Real users upload PDFs and Word files, not Python objects. The extraction agent
takes raw text and produces a validated \texttt{DisputeScenario} in one model
call, with a repair retry if that call fails. Its prompt mirrors the Court
agent's discipline in both directions: do not invent a scoring formula the source
does not contain, and do not summarise away one it does. That second half was
added after a real case lost its published 60/40 weighting during extraction,
which silently downgraded every downstream Court assessment from exact
verification to qualitative review without anyone noticing.

Two guards were added after failures turned up while running the system on
documents I had not hand-picked. A minimum source-length check now raises before
any model call whenever extraction yields under 200 characters, because a scanned
PDF with no text layer had previously produced not an error but an entire
fictional dispute, complete with its own contract value and scoring formula.
Separately, an unstated contract value has to serialise as \texttt{null} and
never as zero, since a zero reads downstream as a genuine \pounds0 contract
rather than as missing information.

\section{Structured output and compliance instrumentation}
\label{sec:compliance-method}

Every agent returns JSON validated against a Pydantic schema. Before this
instrumentation existed, a run in which every single response needed repair
looked identical to one in which none did, so a counter layer was added around
parsing. It records brace-extraction fallbacks when \texttt{json.loads} fails,
field coercions during validation, repetition regenerations, and outright parse
failures. The headline metric is \emph{structural compliance}: the proportion of
structured responses that needed no repair at all. It is scoped per response
rather than summed over events, because a single response can trigger several
coercions at once, and naively summing them produces a rate that climbs past one
and then goes negative. Repetition retries are counted but do not mark a response
dirty, since repetitiveness is a content problem rather than a structural one.

\section{Method selection, and what was rejected}
\label{sec:alternatives}

Table~\ref{tab:alternatives} records the significant choices made against the
alternatives that were considered. Three of these are worth commenting on
directly.

\begin{table}[htbp]
\centering
\footnotesize
\caption[Design alternatives]{Design decisions against the alternatives considered.
Rows marked \emph{plan change} diverge from the original project plan.}
\label{tab:alternatives}
\begin{tabular}{@{}p{2.6cm}p{3.0cm}p{6.4cm}@{}}
\toprule
\textbf{Decision} & \textbf{Alternative} & \textbf{Reasoning} \\
\midrule
Prompt engineering + context injection & Fine-tuning a domain model &
  No labelled corpus of procurement dispute negotiations exists to fine-tune on,
  and building one was not realistic within the time available. Prompts also stay
  auditable, which matters when fabrication is the thing being tested for. \\
Llama 3.1 8B \cite{grattafiori2024llama}, self-hosted \emph{(plan change)} & GPT-4o via API &
  Cost, and the fact that client documents are commercially sensitive. It is
  also the harder test: a discipline that holds on an 8B model is not quietly
  relying on raw capability to paper over a weak prompt. \\
Single-pass extraction & Retrieval-augmented generation &
  RAG solves retrieving the right authority; the failure mode here is inventing
  facts about the case actually in front of the model, which retrieval does not
  touch. Scoped out explicitly rather than left unaddressed. \\
LangGraph state machine & Free-form agent conversation &
  Deterministic turn order and explicit conditional edges are what make the
  ablation in \S\ref{sec:baselines} possible at all: you can delete a node and
  measure what happens. \\
React frontend \emph{(plan change)} & Streamlit &
  Needed to be demonstrated to a client board, and needed to stream a
  negotiation live while it was still in progress. \\
JSON logs \emph{(plan change)} & PostgreSQL &
  Every analysis in Chapter~\ref{ch:evaluation} is a batch read over completed
  runs; a database would have added operational overhead for no analytical
  benefit. \\
Real judgments only \emph{(plan change)} & Synthetic scenarios &
  The two synthetic scenarios were AI-authored and were deleted, along with
  their evaluation batches, once a real corpus existed. A synthetic dispute
  cannot falsify anything. \\
\bottomrule
\end{tabular}
\end{table}

The move from synthetic to real scenarios is the most consequential of these.
The project plan had proposed synthetic dispute scenarios with real cases kept as
optional validation. In practice the synthetic pair had been written with model
assistance, which quietly made the whole test circular: a scenario generated by
a language model and then assessed by a language model tells you nothing
independent about whether either one is right. They were deleted outright,
along with seven batch directories of results derived from them, rather than
kept with a caveat attached, because those results were still being served into
the frontend's analytics page, where no caveat could actually reach a viewer.

The plan had also proposed ablations over temperature and BATNA values. Those
were dropped in favour of the three baselines in Section~\ref{sec:baselines}, on
the judgement that ``does this architecture beat a single prompt'' is the
question an examiner asks first, and it is not a question a temperature sweep
can answer.

\section{Delivery and the industrial deliverable}
\label{sec:risk-screen}

The negotiation simulator is exposed through a FastAPI backend and a
Fusion21-branded React client, with live negotiations streamed over
server-sent events. Inference runs against local Ollama by default and against a
tunnelled University GPU instance whenever one is reachable, detected at startup
and re-establishable from a button in the interface. That last piece exists
because the tunnel is fragile in practice, and a demonstration that needs a
terminal open beside it is not really a demonstration.

The component Fusion21 actually asked for is a different thing entirely. Their
framing was prevention: every other part of this system operates on a dispute
that has already been raised. The pre-award challenge risk screen is a
rule-based function with no model call in it at all, taking the same 16 profile
fields already encoded from Doc 1 and returning a risk band, a normalised score,
and a list of flags. Each flag names the field that triggered it, the Doc 1
category it traces back to, a severity, a confidence tag, a rationale closely
paraphrased from a specific Doc 1 sentence, and a concrete pre-award mitigation.

Three boundaries here are deliberate. It predicts whether a challenge is likely
to be \emph{raised}, not whether the authority would go on to lose one, because
that second question belongs to the Court agent, and blurring the two would
smuggle a merits judgment in through inputs that really only speak to a bidder's
incentive to complain. Every flag carries a confidence tag, because an authority
knows the quality of its own documentation but can only guess at a losing
bidder's legal representation. And the score is a triage ordering device, not a
probability: Doc 1 gives directional relationships rather than weights, so the
severity weighting here is a transparent first pass, and the module's own output
says as much it is not claiming to be calibrated.

\section{Verification discipline as a method}
\label{sec:verification-discipline}

One methodological commitment runs through everything described above, and it is
worth stating outright as a method rather than leaving it as an implied habit. No
number, citation or case fact enters this project without being checked against
a primary source, and that rule was applied to my own write-ups just as strictly
as to model output. It caught real errors in both directions. An earlier draft of
the evaluation claimed the agents had negotiated over a case's 68\%/84\% score
gap; they had not, because the cached scenario never actually contained it, and
the sentence had been checked against the raw case text rather than against what
the pipeline had actually fed the agents. A contract value of roughly
\pounds1.5 billion attributed to one case turned out, on verification, to be
\pounds112.1 million for the specific procurement actually at issue. Both errors
belong to the same class the Court agent is hardened against, just occurring in
the human layer instead, and both are reported here rather than quietly fixed
and left unmentioned.